\documentclass[12pt]{article}
\usepackage[utf8]{inputenc}
\usepackage[T1]{fontenc}
\usepackage{hyperref}
\title{Kitharaka Language Resources for K-12 Curriculum Development}
\author{Compiled by Claude Sonnet 5 (claude-sonnet-5) for Andr\'as Kornai}
\date{}
\begin{document}
\maketitle

\section{Introduction}
Kitharaka (Tharaka; ISO 639-3 thk) is a Bantu language of Kenya. The searches found a small set of mother-tongue literacy materials, a complete Bible translation, and audio and profile pages. They found no dictionary and no grammar.

\section{Primary Educational Resources}
\begin{itemize}
\item \emph{T\^uthoome K\^i\^itharaka: Let's read Kitharaka}, a self-teaching transition primer (Bible Translation and Literacy, Nairobi, 2000) by John Nguu and Albert Kathenya, for readers of another language who want to read Kitharaka: \url{https://www.syncog.ppls.ed.ac.uk/KiitharakaProject/site/K%C3%AE%C3%AEtharaka_Digital_Materials/Tharaka%20Self-teaching%20Transition%20Primer%20(2000)/Self-teaching%20Transition%20Primer%20(2000).pdf}
\item Tharaka Standard 3 students' book (2003, reprinted 2009), mother-tongue curriculum material developed with the Tharaka District Education Office on the KIE syllabus: \url{https://www.syncog.ppls.ed.ac.uk/KiitharakaProject/site/K%C3%AE%C3%AEtharaka_Digital_Materials/Tharaka%20Std%203%20Students%20book%20(2003,%20reprint%209)/Tharaka%20Std%203%20final%20B5.pdf}
\end{itemize}

\section{Linguistic and Reference Materials}
\begin{itemize}
\item Glottolog bibliographic entry for the prereading book \emph{T\^uthoomeni K\^itharaka 1} (BTL, 1995): \url{https://glottolog.org/resource/reference/id/45874}
\item Joshua Project language profile: \url{https://www.joshuaproject.net/index.php/languages/thk}
\end{itemize}
No dictionary or grammar appeared in the search results.

\section{Religious and Community Resources}
\begin{itemize}
\item Scripture Earth listing (text, audio, apps): \url{https://www.scriptureearth.org/00eng.php?iso=thk}
\item Digital Bible Society entry: \url{https://dbs.org/bibles/THKBTL}
\item Joshua Project, Tharaka people in Kenya: \url{https://joshuaproject.net/people_groups/15463}
\end{itemize}
The New Testament was dedicated in 2001 and the complete Bible in 2019, according to the search summaries.

\section{Digital and Audio Resources}
\begin{itemize}
\item Global Recordings Network audio: \url{https://globalrecordings.net/en/language/thk}
\end{itemize}

\section{Curriculum Development Context}
Mother-tongue materials exist at the early primary level (Standard 3) and for transitional literacy. No materials for upper primary or secondary were found. A curriculum would need to build on these, together with the Bible text as a large written corpus, and would need new dictionary, grammar and orthography work.
\end{document}